\section{Diffuse weights and dense prime supports}
\label{an:section}

This section proves three uniform estimates for the raw hinge formulation of
the interval divisor-cover problem. The first allows an arbitrary prime
support when all weights are sufficiently small, together with an extension
allowing one arbitrary exceptional prime weight. The second allows a
controlled number of exceptional weights, at the cost of an explicit
additive defect. The third treats every weight vector on primes below
\(\exp(o(\sqrt{\log M}))\), and gives an asymptotic coefficient-two
integer cover for source sets supported on those primes. These results do
not establish the unrestricted conjecture.

All logarithms are natural unless a base is indicated, and
\((t)_+=\max(t,0)\). Prime sets in this section are finite. For nonnegative
weights \(w=(w_p)_{p\in P}\), write
\[
 X_w(m)=\sum_{p\in P}w_pv_p(m).
\]
If \(I\) consists of \(M\) consecutive positive integers, define
\[
 \Delta_c(M,I,w)
 =\sum_{a=1}^{M}(X_w(a)-c)_+
   -\sum_{b\in I}(X_w(b)-1)_+.
\]
The assertion \(H_c\) is \(\Delta_c(M,I,w)\le0\). Whenever coprimality
counts at real lengths are used, an interval of length \(T\) means a
half-open interval \((x,x+T]\), with its integer points understood and
\(x\ge0\). For a finite prime set \(Q\), put
\[
 Q^{\Pi}=\prod_{q\in Q}q,
 \qquad \vartheta(Q)=\prod_{q\in Q}(1-1/q).
\]
Both empty products equal one. An integer is \(Q\)-smooth if every prime
divisor belongs to \(Q\); the integer one is included.

\subsection{A uniform consequence of the fundamental lemma}

\begin{lemma}[Uniform sifting at a large parameter]
\label{an:sieve}
There is an absolute constant \(C>0\) such that the decreasing function
\begin{equation}
 e(t)=C\bigl((t/2)^{-t/4}+2^{1-t/2}\bigr),\qquad t\ge4,
 \label{an:base-error}
\end{equation}
has the following property. If \(z\ge2\), \(Q\) is any subset of the
primes at most \(z\), and \(J\) is an interval of real length
\(T\ge z^t\), then
\[
 (1-e(t))T\vartheta(Q)
 \le \#\{m\in J:(m,Q^{\Pi})=1\}
 \le (1+e(t))T\vartheta(Q).
\]
In particular, the estimate is uniform in the interval shift and in the
choice of the subset \(Q\), and \(e(t)\to0\) as \(t\to\infty\).
\end{lemma}

\begin{proof}
The fundamental lemma of the combinatorial sieve, in dimension one, gives
\[
 \#\{m\in J:(m,Q^{\Pi})=1\}
 =T\vartheta(Q)\{1+O(u^{-u/2})\}+O(z^u),\qquad u\ge1.
\]
The constants are absolute here. Indeed, the remainder in each
divisibility count is at most one, while the dimension-one product
condition holds with common constants for every subset of the primes.
This is the fundamental-lemma statement and interval application in
\cite{Lagarias}; see also the interval version for subsets of the small
primes in \cite[Theorem~1.4.2]{GranvilleSound}.

Take \(u=t/2\). The elementary product estimate
\[
 \vartheta(Q)
 \ge \prod_{m=2}^{\lfloor z\rfloor}(1-1/m)
 =\frac{1}{\lfloor z\rfloor}\ge\frac1z
\]
shows that the relative contribution of the additive error is at most
\[
 O\!\left(\frac{z^{t/2}}{T\vartheta(Q)}\right)
 =O(z^{1-t/2})=O(2^{1-t/2}).
\]
Increasing a single absolute constant gives the stated function. Both
terms in that function decrease for \(t\ge4\).
\end{proof}

\subsection{A level comparison for diffuse raw weights}

The next proposition works directly with the unbounded valuations
\(v_p\). The final prime-power block of a carrier is allowed to be
partial; this will be essential for controlling large exponents.

\begin{proposition}[Diffuse level comparison]
\label{an:diffuse-level}
Let \(\ell>\eta>0\), let \(0\le w_p\le\eta\), and put
\[
 \mu=\sum_{p\in P}\frac{w_p}{p-1},\qquad
 r=\lfloor\ell/\eta\rfloor\ge8,\qquad
 j_0=\lfloor r/2\rfloor,\qquad \kappa=r-j_0.
\]
Suppose
\begin{equation}
 (1+e(\kappa))\mu+\eta2^{1-j_0}
 \le(\ell-\eta)(1-e(r)).
 \label{an:diffuse-condition}
\end{equation}
Then, for every real \(T>0\), every interval \(J\) of length \(T\), and
every \(s\ge0\),
\begin{equation}
 \#\{a\le T:X_w(a)>s+\ell\}
 \le \#\{b\in J:X_w(b)>s\}.
 \label{an:diffuse-level-ineq}
\end{equation}
Consequently \(H_{1+\ell}\) holds at every integer length and shift.
\end{proposition}

\begin{proof}
Order the prime occurrences from \(P\) in a source integer \(a\) by
increasing prime, retaining multiplicity and omitting primes outside
\(P\). If \(X_w(a)>s+\ell\), let \(d\) be the product of the first
occurrences up to the first crossing of \(s\), and let \(p\) be its
largest prime divisor. Then \(d>1\) and
\begin{equation}
 X_w(d/p)\le s<X_w(d)\le s+\eta.
 \label{an:partial-crossing}
\end{equation}
All \(P\)-primes below \(p\) occur in \(d\) to their full exponents
in \(a\). Further copies of \(p\) may remain in \(a/d\). In particular,
\[
 X_w(a/d)>\ell-\eta.
\]
Every prime occurrence contributes at most \(\eta\), so \(a/d\) has at
least \(r\) occurrences from \(P\), counting multiplicity. All of them
are at least \(p\). Thus \(a/d\ge p^r\), and \(T/d\ge p^r\).

Let \(\mathcal F_s\) consist of all \(P\)-smooth \(d>1\) satisfying
\eqref{an:partial-crossing} and \(dp^r\le T\), where \(p\) is the
largest prime divisor of \(d\). For such a carrier, put
\[
 Q_p=\{q\in P:q<p\},\qquad V_p=\vartheta(Q_p).
\]
Define \(K_d(b)=0\) when \(d\nmid b\), and otherwise put
\[
 K_d(b)=\mathbf1_{(b/d,Q_p^{\Pi})=1}.
\]
For each positive integer \(b\),
\begin{equation}
 \sum_{d\in\mathcal F_s}K_d(b)
 \le \mathbf1_{X_w(b)>s}.
 \label{an:partial-certificate}
\end{equation}
Indeed, positivity forces the factors of \(b\) at all \(P\)-primes
below \(p\) to be exactly those in \(d\). At \(p\), condition
\eqref{an:partial-crossing} specifies the unique first occurrence
crossing \(s\). Hence at most one term is positive, and it forces
\(X_w(b)>s\).

Write \(U=T/d\). The multiples of \(d\) in \(J\), divided by \(d\),
form the integer points of an interval of length \(U\ge p^r\).
Lemma~\ref{an:sieve} and \eqref{an:partial-certificate} imply
\begin{equation}
 \#\{b\in J:X_w(b)>s\}
 \ge \sum_{d\in\mathcal F_s}(1-e(r))UV_p.
 \label{an:diffuse-target}
\end{equation}

For a fixed carrier, enlarge its possible source cofactors to
\[
 \mathcal A_d=\{t\le U:(t,Q_p^{\Pi})=1, X_w(t)>\ell-\eta\}.
\]
Every member still has at least \(r\) occurrences from \(P\), all at
least \(p\). Therefore
\begin{equation}
 |\mathcal A_d|
 \le\frac1{\ell-\eta}\sum_{t\in\mathcal A_d}X_w(t).
 \label{an:cofactor-markov}
\end{equation}
Expand the sum into prime-power divisibility indicators. If \(q\ge p\),
\(j\le j_0\), and some \(t\in\mathcal A_d\) is divisible by \(q^j\),
then \(t/q^j\) has at least \(r-j\) such occurrences. It follows that
\[
 U/q^j\ge p^{r-j}\ge p^{\kappa}.
\]
Terms for which this length inequality fails are zero before enlarging
the cofactor set. For every other term, dividing by \(q^j\) preserves
coprimality to \(Q_p\). Lemma~\ref{an:sieve} therefore bounds the
contribution of all \(j\le j_0\) by
\begin{equation}
 (1+e(\kappa))UV_p
   \sum_{q\in P,\ q\ge p}w_q\sum_{j\ge1}q^{-j}
 \le(1+e(\kappa))UV_p\mu.
 \label{an:low-powers}
\end{equation}

For \(j>j_0\), discard coprimality and the heavy-cofactor restriction.
Their contribution is at most
\[
 U\sum_{q\in P,\ q\ge p}\frac{w_q}{q^{j_0}(q-1)}
 \le U\eta\sum_{m=p}^{\infty}\frac1{m^{j_0}(m-1)}
 \le\frac{U\eta p^{1-j_0}}{p-1}.
\]
The last estimate uses
\(\sum_{m=p}^{\infty}1/[m(m-1)]=1/(p-1)\). Also
\[
 V_p\ge\prod_{m=2}^{p-1}(1-1/m)=\frac1{p-1},
\]
including \(p=2\). Thus the high-power contribution is at most
\begin{equation}
 UV_p\eta2^{1-j_0}.
 \label{an:high-powers}
\end{equation}
Each heavy source is assigned to one carrier, so its total count is at
most \(\sum_d|\mathcal A_d|\). Equations~\eqref{an:cofactor-markov}--\eqref{an:high-powers}, followed
by \eqref{an:diffuse-condition}, show that
\(|\mathcal A_d|\le(1-e(r))UV_p\). Summing over carriers and using
\eqref{an:diffuse-target} proves \eqref{an:diffuse-level-ineq}.
All counts used real interval lengths, so no extra rounding assumption
is needed. Finally, integrating \eqref{an:diffuse-level-ineq} over
\(s\ge1\) proves \(H_{1+\ell}\).
\end{proof}

\subsection{Large weighted means and the diffuse theorem}

\begin{lemma}[Elementary prefix moments]
\label{an:moments}
Suppose \(M\ge2\), \(P\subseteq[2,M]\), and \(0\le w_p\le\eta\).
Let \(\mathbb E_0\) denote averaging over \(1,\ldots,M\), and put
\(\mu=\sum_{p\in P}w_p/(p-1)\). Then
\begin{equation}
 \mu-3\eta\le\mathbb E_0X_w\le\mu,
 \qquad
 \mathbb E_0(X_w-\mu)^2\le8\eta\mu.
 \label{an:moment-bounds}
\end{equation}
If \(0<\delta\le1\), \(\eta\le\delta^2/32\), and
\(\mu\ge1+\delta\), then
\begin{equation}
 \sum_{a=1}^{M}\min(X_w(a),2)\ge M.
 \label{an:capped-mass}
\end{equation}
\end{lemma}

\begin{proof}
For \(e_p=\lfloor\log M/\log p\rfloor\), the rounding loss from
prime powers \(p^j\le M\) is at most \(\eta\): distinct prime powers
are distinct integers, and there are at most \(M-1\) of them. The
omitted tail for each \(p\) is
\[
 \frac{w_p}{p^{e_p}(p-1)}\le\frac{2\eta}{M},
\]
because \(p^{e_p+1}>M\). There are at most \(M\) supporting primes,
so the total tail is at most \(2\eta\). This proves the mean estimate.

Counting multiples gives, for distinct primes \(p,q\),
\[
 \mathbb E_0[v_pv_q]\le\frac1{(p-1)(q-1)},
 \qquad
 \mathbb E_0[v_p^2]\le\sum_{j\ge1}(2j-1)p^{-j}
 =\frac{p+1}{(p-1)^2}.
\]
Consequently
\[
 \mathbb E_0X_w^2
 \le\mu^2+\sum_p\frac{w_p^2p}{(p-1)^2}
 \le\mu^2+2\eta\mu.
\]
Combining this with the lower mean estimate proves the second assertion
in \eqref{an:moment-bounds}.

If \(1+\delta\le\mu\le4\), the map \(x\mapsto\min(x,2)\) is
one-Lipschitz, whence
\[
 \mathbb E_0\min(X_w,2)
 \ge\min(\mu,2)-\sqrt{8\eta\mu}
 \ge1+\delta-\sqrt{32\eta}\ge1.
\]
If \(\mu\ge4\), Chebyshev's inequality gives
\[
 \Pr_0(X_w<2)\le\frac{8\eta\mu}{(\mu-2)^2}\le8\eta,
\]
and therefore
\(\mathbb E_0\min(X_w,2)\ge2(1-8\eta)\ge1\).
This proves \eqref{an:capped-mass}.
\end{proof}

For any nonnegative raw weights, prime-power divisibility counts imply
\(\sum_{b\in I}X_w(b)\ge\sum_{a\le M}X_w(a)\). Thus the elementary
identity \((x-c)_+=x-\min(x,c)\) gives
\begin{equation}
 \Delta_2(M,I,w)
 \le M-\sum_{a=1}^{M}\min(X_w(a),2).
 \label{an:automatic-hinge}
\end{equation}
In particular, \eqref{an:capped-mass} implies \(H_2\). Adding further
nonnegative prime coordinates preserves this sufficient condition.

\begin{theorem}[Uniform diffuse-weight hinge inequality]
\label{an:diffuse-theorem}
For every \(\varepsilon>0\) there exists \(\eta_\varepsilon>0\) such
that \(H_{2+\varepsilon}\) holds at every positive integer length and
every shift for every finite prime support with
\(0\le w_p\le\eta_\varepsilon\).
\end{theorem}

\begin{proof}
It suffices to treat \(0<\varepsilon\le1\), since increasing the prefix
threshold decreases its hinge. Put \(\ell=1+\varepsilon\) and
\(\delta=\varepsilon/2\). Choose \(\eta>0\) sufficiently small that
\(\eta\le\varepsilon^2/128\), \(r=\lfloor\ell/\eta\rfloor\ge8\),
and \eqref{an:diffuse-condition} holds with \(\mu\) replaced by
\(1+\varepsilon/2\). Such a choice exists: as \(\eta\downarrow0\),
the two sides tend respectively to \(1+\varepsilon/2\) and
\(1+\varepsilon\).

The case \(M=1\) is immediate. For \(M\ge2\), remove primes greater
than \(M\) before computing \(\mu\). They contribute nothing to the
prefix, and restoring them only increases the interval hinge. If
\(\mu\le1+\varepsilon/2\), Proposition~\ref{an:diffuse-level} applies.
If \(\mu\ge1+\varepsilon/2\), Lemma~\ref{an:moments} and
\eqref{an:automatic-hinge} give the stronger assertion \(H_2\).
\end{proof}

\begin{corollary}[One arbitrary exceptional prime]
\label{an:one-heavy}
The conclusion of Theorem~\ref{an:diffuse-theorem} remains valid if one
prime \(q\) has an arbitrary nonnegative weight, while every other
weight is at most \(\eta_\varepsilon\). The exceptional weight need
not be at most one in this auxiliary assertion.
\end{corollary}

\begin{proof}
Write \(V\) for the raw load from primes other than \(q\), and first
discard primes greater than \(M\) as above. Let \(\mu_L\) be the mean
parameter of these light primes. If \(\mu_L\ge1+\varepsilon/2\),
Lemma~\ref{an:moments} gives
\(\sum_{a\le M}\min(V(a),2)\ge M\). Adding the \(q\)-coordinate
preserves this inequality, so \eqref{an:automatic-hinge} proves \(H_2\).

Suppose \(\mu_L\le1+\varepsilon/2\), and put
\(\ell=1+\varepsilon\), \(c=1+\ell\). The proof of
Theorem~\ref{an:diffuse-theorem} supplies the level comparison
\eqref{an:diffuse-level-ineq} for \(V\), at every length and shift and
every nonnegative level. It remains available after restricting to any
smaller length; primes above that length may again be removed and then
restored. Also \(V(q^jt)=V(t)\).

First adjoin a nondecreasing profile \(u(j)\in[0,1]\), with \(u(0)=0\).
For \(h=c\) or \(h=1\), the telescoping hinge identity is
\begin{align}
 (V(n)+u(v_q(n))-h)_+
 &= (V(n)-h)_+ \notag\\
 &\quad+\sum_{j\ge1}\mathbf1_{q^j\mid n}
   \int_{u(j-1)}^{u(j)}\mathbf1_{V(n)>h-v}\,dv.
 \label{an:one-heavy-increments}
\end{align}
The base hinges satisfy \(H_c\). In each increment, the target threshold
is \(s=1-v\ge0\) and the source threshold is \(c-v=s+\ell\).
Multiples of \(q^j\) in the prefix rescale to
\(1,\ldots,\lfloor M/q^j\rfloor\); those in \(I\) rescale to a
consecutive interval containing at least that many integers. Apply the
level comparison on a subinterval of this length. If the length is zero,
the prefix term vanishes. Extra target positions are nonnegative.
Thus each prefix increment is bounded by its target increment, proving
\(H_c\) for \(V+u(v_q)\).

For the raw exceptional weight, take
\[
 u(j)=\min(w_qj,1),\qquad R(j)=(w_qj-1)_+.
\]
Pointwise,
\[
 (V+u+R-c)_+\le R+(V+u-c)_+,
 \qquad
 (V+u+R-1)_+=R+(V+u-1)_+.
\]
The identity holds because \(R>0\) forces \(u=1\). The nondecreasing
\(q\)-adic function \(R\) has prefix sum at most its sum over \(I\),
by expanding its nonnegative successive increments and counting
multiples of \(q^j\). This proves the raw assertion.
\end{proof}

\subsection{A controlled number of exceptional weights}

\begin{lemma}[Sifting with arbitrary exceptional primes]
\label{an:exceptional-sieve}
Let \(E\) contain at most \(k\) primes, and let \(Q\) be a set of
primes at most \(z\ge2\), disjoint from \(E\). Put
\begin{equation}
 e_k(t)=3^ke(t/2)+3\cdot4^k2^{1-t/2},\qquad t\ge8.
 \label{an:exceptional-error}
\end{equation}
If an interval \(J\) has real length \(T\ge z^t\), then
\[
 \left|\#\{m\in J:(m,Q^{\Pi}E^{\Pi})=1\}
       -T\vartheta(Q)\vartheta(E)\right|
 \le e_k(t)T\vartheta(Q)\vartheta(E).
\]
The estimate is uniform in the exceptional primes, including primes
larger than \(T\). For each fixed \(k\), \(e_k(t)\) decreases to zero.
\end{lemma}

\begin{proof}
Split \(E=E_{\le z}\cup E_{>z}\), put
\(Q_0=Q\cup E_{\le z}\), \(V_0=\vartheta(Q_0)\), and
\(F=\prod_{p\in E_{>z}}p\). Inclusion--exclusion, with \(\mu(f)\)
denoting the M\"obius function, gives
\[
 \#\{m\in J:(m,Q^{\Pi}E^{\Pi})=1\}
 =\sum_{f\mid F}\mu(f)\#\{m\in J/f:(m,Q_0^{\Pi})=1\}.
\]
For \(f\le T/z^{t/2}\), Lemma~\ref{an:sieve} applies at length
\(T/f\), with relative error \(e(t/2)\). The sum of these errors is
at most \(e(t/2)TV_0\sum_{f\mid F}1/f\). Relative to
\(TV_0\vartheta(E_{>z})\), it is bounded by
\[
 e(t/2)\prod_{p\in E_{>z}}\frac{1+1/p}{1-1/p}
 \le3^ke(t/2).
\]
For each remaining \(f\), its exact count is at most
\(T/f+1\le z^{t/2}+1\), while the magnitude of its main term is at
most \(z^{t/2}\). There are at most \(2^k\) such divisors, so their
combined error is at most \(3\cdot2^kz^{t/2}\). Since
\(V_0\ge1/z\) and \(\vartheta(E_{>z})\ge2^{-k}\), the relative
error is at most
\[
 3\cdot4^kz^{1-t/2}\le3\cdot4^k2^{1-t/2}.
\]
Combining the two bounds proves the lemma.
\end{proof}

\begin{proposition}[Conditional diffuse level comparison]
\label{an:conditional-level}
Let \(E\) contain at most \(k\) primes, and let
\(V(m)=\sum_{p\in L}w_pv_p(m)\), where \(L\cap E=\varnothing\) and
\(0\le w_p\le\eta\). Put
\[
 \mu_L=\sum_{p\in L}\frac{w_p}{p-1},\quad
 r=\lfloor\ell/\eta\rfloor\ge16,\quad
 j_0=\lfloor r/2\rfloor,\quad \kappa=r-j_0,
\]
where \(\ell>\eta\). If
\begin{equation}
 (1+e_k(\kappa))\mu_L+\eta2^{k+1-j_0}
 \le(\ell-\eta)(1-e_k(r)),
 \label{an:conditional-condition}
\end{equation}
then for every real \(T>0\), interval \(J\) of length \(T\), and
\(s\ge0\),
\begin{equation}
 \#\{a\le T:(a,E^{\Pi})=1, V(a)>s+\ell\}
 \le\#\{b\in J:(b,E^{\Pi})=1, V(b)>s\}.
 \label{an:conditional-level-ineq}
\end{equation}
\end{proposition}

\begin{proof}
Use the first crossing in the ordered light-prime occurrences, allowing
a partial last block, as in Proposition~\ref{an:diffuse-level}. Each
carrier \(d>1\), with largest prime \(p\), satisfies
\[
 V(d/p)\le s<V(d)\le s+\eta,\qquad T/d\ge p^r.
\]
Let \(Q_p=\{q\in L:q<p\}\) and \(V_p=\vartheta(Q_p)\). For
\(d\mid b\), use the certificate
\[
 K_d(b)=\mathbf1_{(b/d,E^{\Pi}Q_p^{\Pi})=1},
\]
and put \(K_d(b)=0\) otherwise. A positive term is the unique first
crossing of \(s\) among the light-prime occurrences of \(b\). Since
\(d\) is coprime to \(E^{\Pi}\), positivity also forces
\((b,E^{\Pi})=1\). Thus the sum of the certificates is bounded by the
target indicator in \eqref{an:conditional-level-ineq}.

Write \(U=T/d\). Lemma~\ref{an:exceptional-sieve} gives a target
contribution at least
\((1-e_k(r))UV_p\vartheta(E)\) for each carrier. Enlarge its source
cofactors to those \(t\le U\) coprime to \(E^{\Pi}Q_p^{\Pi}\)
with \(V(t)>\ell-\eta\), and apply the Markov bound
\eqref{an:cofactor-markov}. Such cofactors have at least \(r\) light
prime occurrences, all at least \(p\).

For \(j\le j_0\), a nonzero term divisible by \(q^j\) forces
\(U/q^j\ge p^{r-j}\ge p^{\kappa}\). Division by \(q^j\) preserves
all the relevant coprimality conditions, because \(q\in L\),
\(q\ge p\), and \(L\cap E=\varnothing\). The upper estimate in
Lemma~\ref{an:exceptional-sieve} therefore bounds the total low-power
contribution by
\[
 (1+e_k(\kappa))UV_p\vartheta(E)\mu_L.
\]
The high-power estimate before inserting the density is unchanged:
it is at most \(U\eta p^{1-j_0}/(p-1)\). Now
\(V_p\ge1/(p-1)\) and \(\vartheta(E)\ge2^{-k}\), so this is at most
\[
 UV_p\vartheta(E)\eta2^{k+1-j_0}.
\]
Divide the sum of these bounds by \(\ell-\eta\), use
\eqref{an:conditional-condition}, and sum over carriers. This proves
the proposition, including its real-length formulation.
\end{proof}

\begin{theorem}[Exceptional coordinates with an explicit defect]
\label{an:fixed-core}
For every \(\varepsilon>0\) there is \(c_\varepsilon>0\) such that
the following holds for every integer \(k\ge0\). Suppose
\(0\le w_p\le1\) and there is a subset \(E\subseteq P\) of size at most \(k\)
such that
\[
 w_p\le\frac{c_\varepsilon}{k+1}\qquad(p\notin E).
\]
Then, uniformly in \(M\), \(I\), and the actual exceptional primes,
\begin{align}
 \Delta_{2+\varepsilon}(M,I,w)
 &\le 2^{|E|+1}
   \sum_{\substack{h\le M\\h\ E\text{-smooth}}}(Z(h)-1)_+
   \label{an:exact-core-defect}\\
 &\le2^{k+1}(1+\log_2M)^{k+1},
 \label{an:core-defect}
\end{align}
where \(Z(h)=\sum_{p\in E}w_pv_p(h)\).
\end{theorem}

\begin{proof}
Again it suffices to consider \(0<\varepsilon\le1\). We first make
the small-weight threshold quantitative. For \(u\ge16\),
\[
 e(u/2)\le3C2^{-u/4},
\]
because \(u/4\ge4\) implies
\((u/4)^{-u/8}\le2^{-u/4}\). Hence
\begin{equation}
 e_k(u)\le3C3^k2^{-u/4}+6\cdot4^k2^{-u/2}.
 \label{an:quant-sieve}
\end{equation}
Choose an integer \(R_\varepsilon\ge16\) so large that
\[
 3C2^{-R_\varepsilon/4}+6\cdot2^{-R_\varepsilon/2}
 \le\varepsilon/16.
\]
For \(u\ge16k+R_\varepsilon\), \eqref{an:quant-sieve} then gives
\(e_k(u)\le\varepsilon/16\). Choose \(c_\varepsilon>0\) satisfying
\[
 c_\varepsilon\le\varepsilon^2/128,
 \qquad
 \frac{1+\varepsilon}{c_\varepsilon}
 \ge\max(32,2R_\varepsilon+2).
\]
Set \(\eta=c_\varepsilon/(k+1)\), \(\ell=1+\varepsilon\), and
define \(r,j_0,\kappa\) as in Proposition~\ref{an:conditional-level}.
Then
\[
 r\ge32k+2R_\varepsilon,\qquad
 \kappa\ge16k+R_\varepsilon,\qquad j_0\ge k+1.
\]
In particular \(e_k(r),e_k(\kappa)\le\varepsilon/16\), and the
tail term \(\eta2^{k+1-j_0}\) is at most \(\eta\). If
\(\mu_L\le1+\varepsilon/2\), these estimates verify
\eqref{an:conditional-condition}. Explicitly, using
\(\eta\le\varepsilon/128\), the difference between its resulting
right-hand lower bound and left-hand upper bound is at least
\[
 \frac{46}{128}\varepsilon-\frac{191}{2048}\varepsilon^2>0.
\]

The case \(M=1\) is immediate. Otherwise discard all primes greater
than \(M\) before computing \(\mu_L\) or splitting into cases.
Restoring them only increases the target hinge. Removed exceptional
primes cannot occur in any \(E\)-smooth \(h\le M\), so this leaves
the smooth-part sum unchanged and can only decrease the bound in
\eqref{an:exact-core-defect}. Write
\(X_w=V+Z\), where \(V\) uses the primes outside \(E\).
If \(\mu_L\ge1+\varepsilon/2\), Lemma~\ref{an:moments} and
\eqref{an:automatic-hinge} give \(H_2\) for the full load, which is
stronger than both asserted bounds.

It remains to consider \(\mu_L\le1+\varepsilon/2\). Each positive
integer has a unique decomposition
\[
 a=ht,\qquad h\ E\text{-smooth},\qquad(t,E^{\Pi})=1.
\]
Here \(h\) is its \emph{exact} exceptional-prime part, and
\(X_w(a)=Z(h)+V(t)\). Fix \(s\ge0\). At a fixed \(h\le M\),
the prefix and target counts are over intervals of the same real length
\(M/h\). If \(s-Z(h)\ge0\), apply
Proposition~\ref{an:conditional-level} at that level; this part has
nonpositive level-count defect.

If \(s-Z(h)<0\), every target cofactor coprime to \(E^{\Pi}\) meets
the target threshold. The source count is at most the unrestricted
prefix coprimality count. Finite inclusion--exclusion shows that each
of these unrestricted counts differs from
\((M/h)\vartheta(E)\) by at most \(2^{|E|}\). Their difference is
therefore at most \(2^{|E|+1}\). Target integers with exact
exceptional part \(h>M\) are nonnegative contributions and may be
discarded. Summing gives
\begin{align}
 &\#\{a\le M:X_w(a)>s+\ell\}
   -\#\{b\in I:X_w(b)>s\} \notag\\
 &\hspace{15mm}\le
 2^{|E|+1}\#\{h\le M:h\ E\text{-smooth},\ Z(h)>s\}.
 \label{an:core-level-defect}
\end{align}
Integrating over \(s\ge1\) proves \eqref{an:exact-core-defect}.
Finally, \(Z(h)\le\Omega(h)\le\log_2M\), and there are at most
\((1+\log_2M)^{|E|}\) exceptional smooth parts \(h\le M\), by
bounding each exponent separately. This proves \eqref{an:core-defect}.
\end{proof}

\begin{corollary}[A growing-core exclusion for raw dual witnesses]
\label{an:growing-core}
Fix \(\varepsilon>0\), and consider weights in \([0,1]\). If
\(k=o(\log M/\log\log M)\) and all but
at most \(k\) weights are at most \(c_\varepsilon/(k+1)\), then
\[
 \Delta_{2+\varepsilon}(M,I,w)\le M^{o(1)}=o(\sqrt M).
\]
More specifically, fix \(\delta\in(0,1)\), let \(A\subseteq[2,M]\)
have \(n\) distinct elements, and suppose
\[
 M\le16n^2,
 \qquad k\le(1-\delta)\frac{\log n}{\log\log n}.
\]
For such weight vectors, as \(n\to\infty\),
\begin{equation}
 \sum_{a\in A}X_w(a)-\sum_{b\in I}(X_w(b)-1)_+
 \le(2+\varepsilon)n+n^{1-\delta+o(1)}.
 \label{an:dual-exclusion}
\end{equation}
\end{corollary}

\begin{proof}
The logarithm of the right side of \eqref{an:core-defect} is
\((k+1)\{\log2+\log(1+\log_2M)\}\). This is \(o(\log M)\) in
the first range. In the second range it is at most
\((1-\delta+o(1))\log n\). Finally,
\[
 \sum_{a\in A}X_w(a)
 \le(2+\varepsilon)n+
      \sum_{a=1}^{M}(X_w(a)-2-\varepsilon)_+,
\]
where distinctness is used to pass from \(A\) to the full prefix.
Combine this with Theorem~\ref{an:fixed-core}.
\end{proof}

Corollary~\ref{an:growing-core} excludes a class of raw dual vectors.
It does not assert a per-instance integer-cover bound equal to the raw
LP value plus \(o(n)\). In particular, a core of size comparable to
a positive power of \(n\) is outside its effective range.

\subsection{Dense support up to a subexponential cutoff}

\begin{theorem}[Uniform dense-support hinge inequality]
\label{an:subexp-hinge}
For every \(\varepsilon>0\) there are
\(\gamma_\varepsilon>0\) and \(M_\varepsilon\) such that
\(H_{2+\varepsilon}\) holds for all \(M\ge M_\varepsilon\), all
shifts, and all weights \(0\le w_p\le1\) supported on primes
\[
 p\le Y\le\exp(\gamma_\varepsilon\sqrt{\log M}),\qquad Y\ge2.
\]
One admissible choice is
\begin{equation}
 J_\varepsilon=\left\lceil\frac2{\eta_\varepsilon}\right\rceil,
 \qquad
 \gamma_\varepsilon=\frac1{\sqrt{16(J_\varepsilon+5)}},
 \label{an:gamma}
\end{equation}
where \(\eta_\varepsilon\) is any threshold from
Theorem~\ref{an:diffuse-theorem}. The constant
\(M_\varepsilon\) is then sufficiently large.
\end{theorem}

\begin{proof}
The empty support is trivial. Give each \(P\)-smooth positive integer
\(h\) probability \(\vartheta(P)/h\). These probabilities sum to one
by the convergent finite Euler product. Under this geometric reference
distribution,
\begin{equation}
 \mathbb E X_w=\sum_{p\in P}\frac{w_p}{p-1}\le1+\log Y.
 \label{an:model-mean}
\end{equation}
If some \(w_p\ge\eta_\varepsilon\), put \(j=\lceil2/w_p\rceil\).
Then \(j\le J_\varepsilon\), and \(X_w(h)\ge2\) whenever
\(p^j\mid h\). The difference
\((X_w-1)_+-(X_w-2)_+\) is nonnegative and equals one on this event.
Its reference mean therefore satisfies
\begin{equation}
 G:=\mathbb E\{(X_w-1)_+-(X_w-2)_+\}
 \ge p^{-j}\ge Y^{-J_\varepsilon}.
 \label{an:model-gap}
\end{equation}

We first bound the contribution of large smooth parts. Set
\[
 u=\frac{\log M}{\log Y},\qquad H=\sqrt M,\qquad
 \sigma=\frac1{4\log Y},\qquad
 R=\sum_{\substack{h>H\\h\ P\text{-smooth}}}\frac{X_w(h)}h.
\]
Since \(X_w(h)\le\log h/\log2\), Rankin's inequality gives
\begin{equation}
 R\le\frac{H^{-\sigma}}{\log2}
 \sum_{h\ P\text{-smooth}}\frac{\log h}{h^{1-\sigma}}.
 \label{an:rankin-start}
\end{equation}
For \(p\le Y\), put \(z_p=p^{-1+\sigma}\). Then
\[
 z_p\le e^{1/4}/p,
 \qquad \frac{z_p}{1-z_p}\le\frac4p,
\]
using \(e^{1/4}<4/3\). Thus
\[
 F:=\prod_{p\in P}(1-p^{-1+\sigma})^{-1}\le Y^4,
\]
because \(\log F\le4\sum_{p\le Y}1/p\le4\log Y\).
Differentiation of this finite Euler product yields
\[
 \sum_{h\ P\text{-smooth}}\frac{\log h}{h^{1-\sigma}}
 =F\sum_{p\in P}\frac{(\log p)z_p}{1-z_p}
 \le4Y^4(\log Y)^2.
\]
Here \(\sum_{p\le Y}(\log p)/p\le(\log Y)^2\), by comparison with
all integers. Since \(H^{-\sigma}=e^{-u/8}\),
\begin{equation}
 R\le\frac4{\log2}Y^4(\log Y)^2e^{-u/8}.
 \label{an:rankin-tail}
\end{equation}

Decompose each positive integer uniquely as \(ht\), where \(h\) is
its exact \(P\)-smooth part and \((t,P^{\Pi})=1\). For \(h\le H\),
both prefix and target counts at that smooth part are coprimality
counts in intervals of real length \(M/h\ge\sqrt M=Y^{u/2}\).
By Lemma~\ref{an:sieve}, their relative errors from
\((M/h)\vartheta(P)\) are at most \(e(u/2)\), provided \(u\ge8\).
The prefix contribution from \(h>H\) is at most \(MR\), since there
are at most \(M/h\) prefix multiples of \(h\). Discard the
nonnegative target contribution from these large parts.

The truncated reference gap loses at most \(R\): indeed
\(0\le(X_w-1)_+-(X_w-2)_+\le X_w\) and
\(\vartheta(P)\le1\). The sum of the two reference hinge means is
at most \(2\mathbb E X_w\). Therefore
\begin{equation}
 \frac{\Delta_2(M,I,w)}M
 \le-Y^{-J_\varepsilon}+2R+2e(u/2)(1+\log Y).
 \label{an:subexp-master}
\end{equation}

We check this comparison uniformly over the entire asserted range of
\(Y\). Write \(L=\log M\), \(y=\log Y\), and use the value
\(\gamma=\gamma_\varepsilon\) in \eqref{an:gamma}. Then
\(\log2\le y\le\gamma\sqrt L\), and the ratio of the bound in
\eqref{an:rankin-tail} to \(Y^{-J_\varepsilon}\) is
\[
 \frac4{\log2}y^2
 \exp\left((J_\varepsilon+4)y-\frac{L}{8y}\right).
\]
Both its polynomial factor and its exponent increase with \(y\).
At the upper endpoint the exponent is
\[
 \left((J_\varepsilon+4)\gamma-\frac1{8\gamma}\right)\sqrt L
 =-\gamma(J_\varepsilon+6)\sqrt L.
\]
The ratio consequently tends to zero uniformly.

For \(u\ge16\), the error function satisfies
\(e(u/2)\le3C2^{-u/4}\). The ratio of the last error in
\eqref{an:subexp-master} to \(Y^{-J_\varepsilon}\) is at most
\[
 6C(1+y)\exp\left(J_\varepsilon y-
                        \frac{(\log2)L}{4y}\right).
\]
Its exponent also increases with \(y\), and is a negative constant
times \(\sqrt L\) at the upper endpoint, since
\[
 J_\varepsilon\gamma^2<\frac1{16}<\frac{\log2}{4}.
\]
It too tends to zero uniformly. Thus, for sufficiently large \(M\),
the errors in \eqref{an:subexp-master} are dominated by the gap,
proving \(H_2\) whenever some weight is at least
\(\eta_\varepsilon\). If every weight is at most this threshold,
Theorem~\ref{an:diffuse-theorem} gives \(H_{2+\varepsilon}\) at every
length without a support restriction. The two cases prove the theorem.
\end{proof}

\subsection{The raw boxed dual and an integer-cover consequence}

We include the LP bridge and the required rounding argument to distinguish
the ensuing integer conclusion from the dual-vector exclusions above.
Let \(A\subseteq[2,M]\) be a set of \(n\) distinct sources. Write
\(\tau(A,I)\) for the least size of a subset \(B\subseteq I\) such
that \(\prod_{a\in A}a\) divides \(\prod_{b\in B}b\). Let
\(P_A\) be the set of primes occurring in \(A\). Put
\(D_p=\sum_{a\in A}v_p(a)\). The raw boxed relaxation is
\begin{equation}
 L(A,I)=\min\left\{\sum_{b\in I}x_b:
 0\le x_b\le1,\quad
 \sum_{b\in I}v_p(b)x_b\ge D_p\quad(p\in P_A)\right\}.
 \label{an:raw-lp}
\end{equation}

\begin{lemma}[Raw dual clipping and row-count rounding]
\label{an:lp-bridge}
The program \eqref{an:raw-lp} is feasible, and its value equals
\begin{equation}
 \max_{0\le w_p\le1}
 \left\{\sum_{a\in A}X_w(a)
          -\sum_{b\in I}(X_w(b)-1)_+\right\},
 \label{an:raw-dual}
\end{equation}
where the weights are on \(P_A\). If \(H_c\) holds for all these
weight vectors, then \(L(A,I)\le cn\). Furthermore,
\begin{equation}
 \tau(A,I)\le L(A,I)+|P_A|.
 \label{an:row-rounding}
\end{equation}
\end{lemma}

\begin{proof}
Distinctness gives \(\prod_{a\in A}a\mid M!\), while
\(M!\mid\prod_{b\in I}b\). Thus selecting every interval position is
feasible. Equivalently, \(\sum_{b\in I}v_p(b)\ge D_p\) for every
source prime. Finite LP duality gives the maximum in
\eqref{an:raw-dual} initially over all \(w_p\ge0\).

If some \(w_p>1\), lower it to one while keeping the other weights
fixed. For every column with \(v_p(b)>0\), its hinge is linear
throughout this change, because the final load is still at least one.
The change in the dual objective is
\[
 (w_p-1)\left(\sum_{b\in I}v_p(b)-D_p\right)\ge0.
\]
Columns with \(v_p(b)=0\) are unchanged. Clipping coordinates in turn
therefore proves \eqref{an:raw-dual}.

If \(H_c\) holds, then every vector in \eqref{an:raw-dual} has
objective at most \(cn\), since
\[
 \sum_{a\in A}X_w(a)
 \le cn+\sum_{a\in A}(X_w(a)-c)_+
 \le cn+\sum_{a=1}^{M}(X_w(a)-c)_+.
\]
For rounding, choose an optimal vertex of the feasible boxed polytope.
It has at most \(|P_A|\) coordinates strictly between zero and one.
Otherwise a nonzero perturbation supported on the fractional
coordinates could annihilate all prime rows, and sufficiently small
perturbations in both signs would preserve feasibility and the box,
contradicting extremality. Round each fractional coordinate upward to
one. The nonnegative valuation rows remain satisfied, and the cost
increases by at most \(|P_A|\), proving \eqref{an:row-rounding}.
\end{proof}

\begin{corollary}[A coefficient-two cover for subexponentially smooth sources]
\label{an:smooth-cover}
Let \(Y(M)\ge2\) be a fixed function satisfying
\[
 \log Y(M)=o(\sqrt{\log M}).
\]
Let \(A\) range over sets of \(n\) distinct integers at least two,
with \(M=\max A\), such that every source prime is at most \(Y(M)\).
Then, as \(n\to\infty\), every positive \(M\)-term interval admits
a product cover with
\[
 \tau(A,I)\le(2+o(1))n,
\]
uniformly in the source set subject to this condition and in the interval
shift. In particular this applies to
\(Y(M)=\exp((\log M)^a)\) for every fixed \(0<a<1/2\).
\end{corollary}

\begin{proof}
Fix \(\varepsilon>0\). For all sufficiently large \(M\), the cutoff
is below \(\exp(\gamma_\varepsilon\sqrt{\log M})\).
Theorem~\ref{an:subexp-hinge} and Lemma~\ref{an:lp-bridge} give
\[
 \tau(A,I)\le(2+\varepsilon)n+\pi(Y(M)).
\]
If \(M<16n^2\), then \(n>\sqrt M/4\), whereas
\(\pi(Y(M))\le Y(M)=M^{o(1)}=o(\sqrt M)\). The rounding term is
therefore \(o(n)\) in this range.

If \(M\ge16n^2\), the quadratic coefficient-two theorem of
Wang~\cite{Wang} gives the stronger exact estimate
\(\tau(A,I)\le2n\). Its hypotheses are precisely a set of \(n\)
distinct positive sources at least two, ambient length \(M=\max A\),
and an arbitrary positive interval of that length. Finally,
\(M\ge n+1\) ensures \(M\to\infty\) whenever \(n\to\infty\).
Since \(\varepsilon\) was arbitrary, the two ranges prove the
corollary.
\end{proof}

The cutoff constant \(\gamma_\varepsilon\) depends on
\(\varepsilon\). The preceding corollary therefore uses
\(\log Y=o(\sqrt{\log M})\); it does not assert coefficient two for
one fixed positive multiple of \(\sqrt{\log M}\) in the exponent at
all accuracies. The conclusion also relies on the explicit small
prime-row count in \eqref{an:row-rounding}, rather than a general
fractional-to-integer error assumption. Polynomial prime scales
\(M^\alpha\), and their interactions with a dense small-prime part,
remain outside the support theorem proved here.
